\documentclass[a5paper, 10pt]{article}


\usepackage{cite, amsmath, amssymb, booktabs, lastpage, graphicx, float}
\usepackage[hidelinks]{hyperref}

\usepackage{geometry}
\geometry{margin=1.5cm,a5paper,twoside,inner=2cm}

\usepackage{setspace}


\usepackage{amsthm}

\theoremstyle{plain}
\newtheorem{lemma}{Lemma}
\newtheorem{theorem}{Theorem}
\newtheorem{conjecture}{Conjecture}
\newtheorem{proposition}[theorem]{Proposition}
\newtheorem*{corollary}{Corollary}
\renewcommand{\qedsymbol}{\rule{.7em}{.7em}}

\theoremstyle{remark}
\newtheorem*{remark}{Remark}
\newtheorem*{note}{Note}
\newtheorem{case}{Case}

\theoremstyle{definition}
\newtheorem{definition}{Definition}
\newtheorem{example}{Example}

%% Changing the title, author, etc.
\usepackage{graphicx}
\makeatletter
\renewcommand{\maketitle}{
    \begin{center}
    \rule[.2em]{\textwidth}{0.353mm}
        \begin{minipage}[m]{0.35\textwidth}
            {\scriptsize
                \begin{center}
                    \textsf{\textbf{\huge MATCH}\\
                        \textit{Communications in Mathematical\\
                            and in Computer Chemistry}
                    }
            \end{center}}
        \end{minipage}\hfill
        \begin{minipage}[m]{0.65\textwidth}
            \begin{flushright}
                \baselineskip=10px
                {\scriptsize\sffamily{\itshape  MATCH Commun. Math. Comput. Chem.}
                \textbf{\vol{}}
                (\pubyear)
                    \thepage--\pageref{LastPage}}\\
                {\scriptsize\sffamily {\bfseries ISSN:} 0340--6253}\\
                {\scriptsize\sffamily \textbf{doi:} \doi{10.46793/match.\vol-\num.\id}}
            \end{flushright}
        \end{minipage}
        \rule[1em]{\textwidth}{.353mm}
        \baselineskip=0.30in
        {\Large\bfseries \@title} \par
        \vspace{5mm}
        \baselineskip=0.2in
        {\large\bfseries \@author}\par
        \vspace{1mm}
        {\it \@address} \par
        {\small\tt \@email} \par
        \vspace{3mm}
        {\small (Received \@date)} \par
    \end{center}
    \vspace{3mm}
}
\newcommand{\address}[1]{\def\@address{#1}}
\newcommand{\email}[1]{\def\@email{#1}}
\newcommand{\doi}[1]{\href{https://doi.org/#1}{#1}}

\newcommand{\vol}{00}
\newcommand{\num}{00}
\newcommand{\pubyear}{0000}
\newcommand{\id}{xxxxx}
%%setting page-number
\setcounter{page}{1}


\thispagestyle{empty}
\makeatother

%% Making << \acknowledgment >> command
\newcommand{\acknowledgment}[1]{\vspace{5mm}\singlespacing
    {\noindent\textbf{\textit{Acknowledgment\/}:} #1}
}

%% Setting page numbers
\usepackage{fancyhdr}

%% now reset the headers and footers
\fancyhead{}
\fancyfoot{}
\renewcommand{\headrulewidth}{.3pt}
\voffset = 18pt
\headsep = 3pt

\renewcommand*{\thefootnote}{\fnsymbol{footnote}}

\usepackage[margin=1cm,%
            font=footnotesize,%
            format=hang,%
            labelsep=period,%
            labelfont=bf]{caption}

%%%%%%%%%%%%%%%%%%%%%%%%%%%%%%%%%%%%%%%%%%%%%%%%%%
% PLACE ADDITIONAL PACKAGES HERE (if it is needed):
\usepackage[T1]{fontenc}
\usepackage{lmodern}
\usepackage{mathtools}
\usepackage{array}
\usepackage{enumitem}
\setlist[itemize]{topsep=2pt,itemsep=2pt}
\setlist[enumerate]{topsep=2pt,itemsep=2pt}
\emergencystretch=3em
\hbadness=2000
\setlength{\marginparwidth}{2cm}
%%%%%%%%%%%%%%%%%%%%%%%%%%%%%%%%%%%%%%%%%%%%%%%%%%

%% make the odd pages have the section name on the top right
\fancyhead[RO]{\footnotesize \thepage}
%% make the even pages have the chapter name on the top left
\fancyhead[LE]{\footnotesize \thepage}

%% making page-numbers visible
\pagestyle{fancy}

% Custom macros for the Loyola index family
    \newcommand{\Wuzi}{LO(G;\alpha,\beta,\gamma)}
  \newcommand{\edgepsi}{\psi(x,y;\alpha,\beta,\gamma)}
  \newcommand{\bid}{\mathcal{F}_{\mathrm{BID}}}
% Path to figures
\graphicspath{{figures/}{./}{../figures/}}

%%%%%%%%%%%%%%%%%%%%%%%%%%%%%%%%%%%%%%%%%%%%%%%%%%%%%%%%%%%%%%%%%%%%%%%%%%%%%%%%%%%%%%%

\title{The Loyola Index Family, Extremal Graph Characterizations and Sensitivity Analysis}
\author{Advik Natarajan$^{a}$\footnote{Corresponding author.}, Ganesh Shiwakoti$^{b,}$, Micheal Arockiaraj$^{a}$}

\address{$^a$Department of Mathematics, Loyola College, Chennai, India\\
    $^b$Department of Mathematics and Statistics, Florida Atlantic University, USA
    }

\email{adviknatarajan@gmail.com}

\date{\today}

% !TeX program = pdflatex

\begin{document}

\maketitle

\begin{abstract}
  Several degree-based topological indices
  have been proposed since the 1970s, including the Zagreb,
  Randi\'c, Geometric-Arithmetic, Sum-Connectivity, Harmonic, and Sombor
  indices. In this paper, we
  introduce the \emph{Loyola index family}, a three-parameter
  degree-based topological index built from
  the product, the sum, and the normalized difference of each pair
  of adjacent degrees. For
  $\alpha, \beta, \gamma \in \mathbb{R}$, the \emph{Loyola index}
  is defined as,
  \begin{equation}
  LO(G;\alpha,\beta,\gamma)
  \;=\; \sum_{uv \in E(G)}
  (d_u d_v)^{\alpha}\,(d_u + d_v)^{\beta}\,
  \exp\!\left(\gamma\,\frac{|d_u - d_v|}{d_u + d_v}\right)
  \end{equation} At specific values of its three parameters,
  the \emph{Loyola index family} reduces to many of the classical
  degree-based topological indices.
  Closed-form values of the \emph{Loyola index} for standard graph classes
  are derived. Some sharp upper and lower bounds are also derived.
  Sensitivity analysis is carried out on octane isomers, non-isomorphic
  trees, decane isomers, and the correlations of the \emph{Loyola index}
  are compared with several classical topological indices.
  \end{abstract}

\onehalfspacing

% =====================================================================
\section{Introduction}\label{sec:intro}
  % =====================================================================

  The literature on degree-based topological indices in chemical
  graph theory has grown substantially since the seminal Wiener
  index, the Randi\'c index,
  and the Zagreb indices of Gutman and
  Trinajsti\'c.
  Among the many parametric extensions are the general Randi\'c
  index $R_{\alpha}$~\cite{BollobasErdos1998}, the general
  sum-connectivity index
  $\chi_{\alpha}$~\cite{ZhouTrinajstic2009sumconn}, the forgotten
  index $F$ revisited by Furtula and
  Gutman~\cite{FurtulaGutman2015F}, the geometric--arithmetic index
  $GA$~\cite{VukicevicFurtula2009GA}, the Albertson irregularity
  index~\cite{Albertson1997} and the inverse $\sigma$-index
  problem~\cite{GutmanToganYurttasCevikCangul2018Sigma}.

  While historical literature focused primarily on algebraic formulations
  to capture molecular branching and structural invariants, a distinct
  shift toward geometric interpretations occurred when Gutman introduced
  the Sombor index family \cite{Gutman2021Sombor, Gutman2021SomborBasic}.
  This formulation led to work such as establishing fundamental bounds
  \cite{MilovanovicMilovanovicMatejic2021Sombor} and evaluating how these
  new geometric invariants correlated with traditional degree-based metrics
  \cite{WangMaoLiFurtula2022SomborRelations}. Much of this initial work
  concentrated on identifying extremal bounds within fundamental graph
  classes, particularly optimizing the general and standard variants of the
  index over molecular trees and related networks
  \cite{DasCevikCangulShang2021Sombor, DasGutman2022SomborTrees,
  AhmadFarooqDas2025Extremal}.

   Contemporary studies have introduced modified metrics such as the mean
  Sombor index \cite{MendezBermudezAguilarSanchez2022MeanSombor}, the
  elliptic Sombor index \cite{GutmanFurtulaOz2024Elliptic}, and the
  hyperbolic Sombor index \cite{BarmanDas2026HSO}. Concurrently, parameters
  like the diminished Sombor index \cite{Movahedi2025DSO,
  MovahediGutmanRedzepovicFurtula2026} and the hybrid vertex-edge KG-Sombor
  index \cite{KosariDehgardiKhan2023Lower} have been developed to handle
  unique structural constraints. Beyond purely theoretical analysis, these
  descriptors are increasingly integrated into practical quantitative
  structure-property relationship (QSPR) models and spectral evaluations.
  For instance, hybrid geometric-harmonic-Zagreb descriptors are now
  utilized to optimize the predictive analysis of polycyclic aromatic
  hydrocarbons \cite{ArockiarajPaulClement2023Hybrid}, while modified
  reverse degree invariants help characterize the graph energy and entropy
  of complex 2D metal-organic frameworks \cite{KalaamGreeniArockiaraj2024}.
  This rapid diversification of structural metrics is currently being
  unified through generalized matrix methods
  \cite{ZuoRatherImranAli2022Modified}, dual coindex formulations
  \cite{GujiWali2025Lanzhou}, simplicial network assessments
  \cite{Shang2022Simplicial}, and finite-dimensional linear frameworks
  \cite{Rada2026LinearVDB}, as documented in recent comprehensive
  literature surveys \cite{Gutman2026SomborSurvey}.

  In this paper, we introduce
  the \emph{Loyola index family} as a controlled three-parameter
  Bond-Incident Degree (BID) index. The bounds of
  Section~\ref{sec:bounds} locate this behaviour inside the
  standard BID framework. The exact reductions of the \emph{Loyola index}
  show that this topological index recovers many of the
  classical indices.



  Direct substitution into Equation~\eqref{eq:wuzi} shows that the
  Loyola family contains a wide range of classical BID indices as
  exact reductions at specific parameter triples, either
  identically or up to a multiplicative scaling constant. The reductions
  verified by direct substitution
  include the edge count $m$ at $(0, 0, 0)$, the first Zagreb
  index $M_1$ at $(0, 1, 0)$, the second Zagreb
  index $M_2$ at $(1, 0, 0)$, the hyper-Zagreb
  index $HM$ at $(0, 2, 0)$, the modified second Zagreb index
  ${}^{m}\!M_2$ at
  $(-1, 0, 0)$, the Randi\'c index $R$ at
  $(-1/2, 0, 0)$, the sum-connectivity index
  $\chi$ at $(0, -1/2, 0)$, the
  half harmonic index $H/2$ at
  $(0, -1, 0)$, the inverse sum indeg index
  $ISI$ at $(1, -1, 0)$, the
  geometric-arithmetic index
  $GA$ (up to a factor of $2$) at
  $(1/2, -1, 0)$, and the arithmetic-geometric index $AG$ (up to a
  factor of $2$) at $(-1/2, 1, 0)$. The whole general Randi\'c
  family $R_\alpha$~\cite{BollobasErdos1998} is recovered as the
  $\beta = \gamma = 0$ axis $LO(G; \alpha, 0, 0)$, and the general
  sum-connectivity family
  $\chi_\beta$~\cite{ZhouTrinajstic2009sumconn} as the
  $\alpha = \gamma = 0$ axis $LO(G; 0, \beta, 0)$. The Loyola family
  thus interpolates among the major degree-based indices.


  On the mathematical side, this paper comprises the
  closed-form values on standard graph classes, sharp upper and
  lower bounds in terms of the basic graph parameters and of other
  classical indices. On the
  chemometric side, by performing a sensitivity analysis  similar
  to~\cite{MovahediGutmanRedzepovicFurtula2026}, we observe that
   the Loyola index
  family  falls in the same correlation regime as the classical BID
  indices rather than above them.

  The distinguishing feature of the Loyola family within the
  degree-based BID class is the multiplicative
  $\exp(\gamma \cdot \mathrm{imb}(d_u, d_v))$ factor, where
  $\mathrm{imb}(x, y) = |x - y|/(x + y) \in [0, 1)$ is the
  normalized degree imbalance. The same imbalance ratio appears
  in the diminished Sombor
  index~\cite{MovahediGutmanRedzepovicFurtula2026} and in
  irregularity-based BID variants; the Loyola family combines it
  multiplicatively with the Randi\'c-style $(d_u d_v)^\alpha$ and
  sum-connectivity-style $(d_u + d_v)^\beta$ factors in a single
  three-parameter object. The factor is bounded above by
  $e^{\gamma}$ and reduces to $1$ on regular graphs.

  Also, it is useful to demonstrate the descriptor-space ceiling in
  which the entire family lives. On hydrogen-suppressed molecular
  graphs of maximum degree at most $4$, every BID index is a linear
  functional of the ten unordered edge-degree-pair counts
  $m_{ij}(G) = |\{uv \in E(G) : \{d_u, d_v\} = \{i, j\}\}|$ for
  $1 \le i \le j \le 4$; this finite-dimensional view of
  vertex-degree-based indices has recently been placed on a formal
  linear-algebraic footing by Rada~\cite{Rada2026LinearVDB}. Thus
  the Loyola family, like every BID family on this molecular class,
  lies inside a fixed ten-dimensional descriptor
  space, no matter how
  its three real parameters are chosen; we make this precise in
  Theorem~\ref{thm:saturation}.

% =====================================================================
\section{Preliminaries}\label{sec:prelim}
% =====================================================================

\subsection{Notation}\label{sec:notation}

Let $G = (V, E)$ be a finite simple connected graph of order
$n = |V|$ and size $m = |E|$. For a vertex $v \in V$, the degree
$d_v = d_G(v)$ counts the edges incident to $v$. The maximum and
minimum degrees are denoted $\Delta = \Delta(G)$ and
$\delta = \delta(G)$. In the chemometric sections we work with
hydrogen-suppressed molecular graphs, where the assumption
$\Delta \le 4$ is natural for organic molecules built from
$\{$C, N, O, S, P, halogens$\}$.

The classical graph families used throughout are: the complete
graph $K_n$, the cycle $C_n$, the path $P_n$, the star $S_n$, the complete
bipartite graph $K_{p,q}$, the hypercube $Q_k$ (the
$k$-dimensional binary hypercube; $2^k$ vertices,
$k \cdot 2^{k - 1}$ edges, $k$-regular), the wheel $W_n$ (an
$(n - 1)$-cycle plus a universal vertex), the friendship graph
$F_p$ ($p$ triangles sharing a common vertex), and the
$r$-regular graphs.

For an edge $e = uv \in E(G)$ we write $(d_u, d_v)$ for its
\emph{degree pair} understood as an unordered multiset
$\{d_u, d_v\}$.

\subsection{Definition}\label{sec:def}

\begin{definition}[Loyola index family]\label{def:wuzi}
Let $G = (V, E)$ be a simple connected graph. For
$\alpha, \beta, \gamma \in \mathbb{R}$, the \emph{Loyola index family}
is defined as,
\begin{equation}\label{eq:wuzi}
\Wuzi
\;=\; \sum_{uv \in E(G)}
(d_u d_v)^{\alpha}\,(d_u + d_v)^{\beta}\,
\exp\!\left(\gamma\,\frac{|d_u - d_v|}{d_u + d_v}\right)
\end{equation}
\end{definition}

Equation~\eqref{eq:wuzi} is well defined for all
$\alpha, \beta \in \mathbb{R}$ on graphs without isolated vertices
(so that $d_u, d_v \ge 1$ on every edge). The expression is
invariant under permutations of $(u, v)$ because
$d_u d_v = d_v d_u$, $d_u + d_v = d_v + d_u$, and
$|d_u - d_v| = |d_v - d_u|$; hence $\Wuzi$ depends only on the
multiset of unordered edge degree pairs of~$G$. When $\gamma = 0$
the family reduces to the standard two-parameter
$(\alpha, \beta)$ degree-based family; the exact reductions to
classical BID indices at specific parameter triples are
catalogued in Section~\ref{sec:intro}.
\subsection{Closed-form values on standard graphs}\label{sec:closedforms}

  \begin{proposition}\label{prop:closedforms}
  Let $\alpha, \beta, \gamma \in \mathbb{R}$. Then:
  \begin{enumerate}
  \item  $LO(K_n; \alpha, \beta, \gamma) =
    \dfrac{n(n-1)}{2}\,2^{\beta}(n - 1)^{2\alpha + \beta}$ for $n \ge 2$.
  \item  $LO(C_n; \alpha, \beta, \gamma) = n \cdot 4^{\alpha +
  \beta}$
    for $n \ge 3$.
  \item
    $LO(P_n; \alpha, \beta, \gamma) =
    \begin{cases}
      2^{\beta}, & n = 2,\\[2pt]
      2 \cdot 2^{\alpha} 3^{\beta} e^{\gamma/3}, & n = 3,\\[2pt]
      2 \cdot 2^{\alpha} 3^{\beta} e^{\gamma/3} + (n - 3) \cdot 4^{\alpha +
  \beta},
         & n \ge 4.
    \end{cases}$
  \item  $LO(S_n; \alpha, \beta, \gamma) =
    (n - 1)\,(n - 1)^{\alpha} n^{\beta} \exp\!\left(\gamma\,\dfrac{n -
  2}{n}\right)$
    for $n \ge 2$.
  \item  $LO(K_{p,q}; \alpha, \beta, \gamma) =
    pq\,(pq)^{\alpha}(p + q)^{\beta}\exp\!\left(\gamma\,\dfrac{|p - q|}{p +
  q}\right)$.
  \item  $LO(Q_k; \alpha, \beta, \gamma) =
    k \cdot 2^{k - 1} \cdot 2^{\beta} k^{2\alpha + \beta}$ for $k \ge 1$.
  \item  For $n \ge 5$, $W_n$ has $n - 1$ spoke edges with
    degree pair $(3, n - 1)$ and $n - 1$ rim edges with degree pair
    $(3, 3)$, so
    \[
    LO(W_n; \alpha, \beta, \gamma)
    = (n - 1)\Bigl[\,3^{\alpha} (n - 1)^{\alpha}(n + 2)^{\beta}
         \exp\!\left(\gamma\,\tfrac{n - 4}{n + 2}\right)
       + 9^{\alpha} 6^{\beta}\Bigr].
    \]
  \item  For $p \ge 1$, $F_p$ has $2p$ spoke edges
    $(2, 2p)$ and $p$ outer edges $(2, 2)$, so
    \[
    LO(F_p; \alpha, \beta, \gamma)
    = p \cdot 4^{\alpha + \beta}
     + 2p\,(4p)^{\alpha}(2p + 2)^{\beta}
         \exp\!\left(\gamma\,\tfrac{p - 1}{p + 1}\right).
    \]
  \item  For an $r$-regular graph $G$ with $m$ edges,
    $\Wuzi = m \cdot 2^{\beta} r^{2\alpha + \beta}$.
  \end{enumerate}
  \end{proposition}

  \begin{proof}
  Each case is obtained by enumerating the edge degree pairs of the
  graph and applying Definition~\ref{def:wuzi}. For $r$-regular
  graphs and for $K_n, C_n, Q_k$ (all regular), every edge has
  $d_u = d_v = r$, so $\mathrm{imb} = 0$ and
  $(d_u d_v)^{\alpha}(d_u + d_v)^{\beta}
     = r^{2\alpha}(2r)^{\beta} = 2^{\beta} r^{2\alpha + \beta}$.
  Items (1), (2), (6), and (9) follow immediately. For $S_n$ all
  $n - 1$ edges have degree pair $(1, n - 1)$ with
  $\mathrm{imb} = (n - 2)/n$; item (4) follows. For $K_{p, q}$ all
  $pq$ edges have degree pair $(p, q)$; item (5) follows. For $P_n$
  ($n \ge 4$) there are exactly two pendant edges with pair $(1, 2)$
  and $n - 3$ interior edges with pair $(2, 2)$; item (3) follows by
  direct summation. The boundary cases $n = 2, 3$ are immediate.
  Items (7) and (8) are obtained by partitioning the edges of $W_n$
  and $F_p$ into their two edge-degree-pair classes respectively and
  summing.
  \end{proof}

 % =====================================================================
  \section{Bounds for the Loyola family}\label{sec:bounds}
  % =====================================================================

  This section establishes a coherent set of upper and lower bounds
  on $\Wuzi$. The development proceeds from a bracket bound via the
  edge-contribution function $\psi$,  bounds in $(n, m, \Delta, \delta)$
  for the non-negative
  parameter region, Cauchy--Schwarz and
  Jensen-type bounds in the first two Zagreb indices $M_1$ and
  $M_2$,
  a uniform ratio-bound theorem, and bounds via the generalised
  Randi\'c, sum-connectivity, and Sombor indices. Throughout, $G$
  is a simple connected graph with $\delta \le d_v \le \Delta$ for
  every $v \in V(G)$.

  The summand in~\eqref{eq:wuzi} depends only on the unordered
  edge degree pair $(d_u, d_v)$. We isolate this dependence as follows:

  \begin{definition}\label{def:psi}
  For real parameters $\alpha, \beta, \gamma$ and positive reals
  $x, y$,
  \begin{equation}\label{eq:psi}
  \edgepsi = (xy)^{\alpha}(x + y)^{\beta}
  \exp\!\left(\gamma\,\tfrac{|x - y|}{x + y}\right).
  \end{equation}
  \end{definition}

  With this notation $\Wuzi = \sum_{uv \in E(G)} \psi(d_u, d_v; \alpha,
  \beta, \gamma)$.

  \begin{proposition}\label{thm:bracket}
  Let $G$ be a simple connected graph of size $m \ge 1$ with
  degrees in $[\delta, \Delta]$. Then
  \begin{equation}\label{eq:bracket}
  m \cdot \psi_{\min}(\delta, \Delta; \alpha, \beta, \gamma)
  \;\le\; \Wuzi \;\le\;
  m \cdot \psi_{\max}(\delta, \Delta; \alpha, \beta, \gamma),
  \end{equation}
  where,
  $\psi_{\min} = \min\{\psi(i, j; \alpha, \beta, \gamma)
                : \delta \le i \le j \le \Delta\}$
  and $\psi_{\max} = \max\{\psi(i, j; \alpha, \beta, \gamma)
                : \delta \le i \le j \le \Delta\}$. The upper
  inequality is an equality if and only if every edge of $G$ has a
  degree pair attaining $\psi_{\max}$, and the lower
  inequality is an equality if and only if every edge has a degree
  pair attaining $\psi_{\min}$.
  \end{proposition}

  \begin{proof}
  Because $G$ has $\delta \le d_v \le \Delta$ for every
  $v \in V(G)$, every edge $uv \in E(G)$ has its degree
  pair $\{d_u, d_v\}$ in the admissible set
  $\mathcal{P}(\delta, \Delta) := \{(i, j) : \delta \le i \le j \le
  \Delta\}$.
  By the definitions of $\psi_{\min}$ and $\psi_{\max}$ as the
  minimum and maximum of $\psi$ over $\mathcal{P}(\delta, \Delta)$,
  \[
  \psi_{\min}(\delta, \Delta; \alpha, \beta, \gamma)
  \;\le\;
  \psi(d_u, d_v; \alpha, \beta, \gamma)
  \;\le\;
  \psi_{\max}(\delta, \Delta; \alpha, \beta, \gamma)
  \qquad. \]
  Summing each of the two per-edge inequalities over the $m$ edges
  of $G$ and using $\sum_{uv \in E(G)} \psi(d_u, d_v; \alpha, \beta,
  \gamma) = \Wuzi$
  gives
  \[
  m \cdot \psi_{\min}(\delta, \Delta; \alpha, \beta, \gamma)
  \;\le\;
  \Wuzi
  \;\le\;
  m \cdot \psi_{\max}(\delta, \Delta; \alpha, \beta, \gamma),
  \]

  For the equality cases, the upper inequality
  $\Wuzi \le m \cdot \psi_{\max}$ is the sum of $m$ per-edge
  inequalities $\psi(d_u, d_v; \alpha, \beta, \gamma) \le \psi_{\max}$;
  this sum is an equality if and only if every per-edge inequality
  is, equivalently every edge $uv$ has a degree pair attaining
  $\psi_{\max}$ over $\mathcal{P}(\delta, \Delta)$. The lower
  equality case is the dual.
  \end{proof}

  \begin{theorem}\label{thm:cs}
  Let $\alpha_1 + \alpha_2 = 2\alpha$, $\beta_1 + \beta_2 = 2\beta$,
  and $\gamma_1 + \gamma_2 = 2\gamma$. Then for every simple
  connected graph $G$ without isolated vertices,
  \[
  LO(G; \alpha, \beta, \gamma)
  \;\le\;
  \sqrt{LO(G; \alpha_1, \beta_1, \gamma_1)
         \cdot LO(G; \alpha_2, \beta_2, \gamma_2)}
  \]
  Equality holds if and only if the ratio
  $\psi(d_u, d_v; \alpha_1, \beta_1, \gamma_1)/\psi(d_u, d_v; \alpha_2,
  \beta_2, \gamma_2)$
  is the same for every edge of $G$.
  \end{theorem}

  \begin{proof}
  For each edge $e = uv \in E(G)$ and each $k \in \{1, 2\}$ write
  \[
  \psi_e^{(k)}
  \;:=\; \psi(d_u, d_v; \alpha_k, \beta_k, \gamma_k)
  \;=\; (d_u d_v)^{\alpha_k}\,(d_u + d_v)^{\beta_k}\,
  \exp\!\left(\gamma_k\, \tfrac{|d_u - d_v|}{d_u + d_v}\right).
  \]
  Both $\psi_e^{(1)}$ and $\psi_e^{(2)}$ are strictly positive on
  every edge of a graph without isolated vertices.

  By the additivity of the exponents and the hypothesis
  $\alpha_1 + \alpha_2 = 2\alpha$,
  $\beta_1 + \beta_2 = 2\beta$, $\gamma_1 + \gamma_2 = 2\gamma$,
  \begin{align*}
  \psi_e^{(1)}\,\psi_e^{(2)}
  &= (d_u d_v)^{\alpha_1 + \alpha_2}\,(d_u + d_v)^{\beta_1 + \beta_2}\,
     \exp\!\left((\gamma_1 + \gamma_2)\, \tfrac{|d_u - d_v|}{d_u +
  d_v}\right) \\
  &= (d_u d_v)^{2\alpha}\,(d_u + d_v)^{2\beta}\,
     \exp\!\left(2\gamma\, \tfrac{|d_u - d_v|}{d_u + d_v}\right) \\
  &= \psi(d_u, d_v; \alpha, \beta, \gamma)^{2}.
  \end{align*}
  Taking square roots gives
  $\sqrt{\psi_e^{(1)}\,\psi_e^{(2)}} = \psi(d_u, d_v; \alpha, \beta,
  \gamma)$
  for every edge. Apply the Cauchy--Schwarz inequality to the
  positive sequences $a_e = \sqrt{\psi_e^{(1)}}$,
  $b_e = \sqrt{\psi_e^{(2)}}$ indexed over $E(G)$:
  \[
  \sum_{e \in E(G)} a_e\, b_e
  \;\le\;
  \sqrt{\sum_{e \in E(G)} a_e^{2}}\,
  \sqrt{\sum_{e \in E(G)} b_e^{2}}.
  \]
  Substituting $a_e b_e = \sqrt{\psi_e^{(1)} \psi_e^{(2)}}
     = \psi(d_u, d_v; \alpha, \beta, \gamma)$,
  $\sum_e a_e^2 = \sum_e \psi_e^{(1)} = LO(G; \alpha_1, \beta_1,
  \gamma_1)$,
  and $\sum_e b_e^2 = \sum_e \psi_e^{(2)} = LO(G; \alpha_2, \beta_2,
  \gamma_2)$
  gives
  \[
  \Wuzi
  \;=\; \sum_{e \in E(G)} \psi(d_u, d_v; \alpha, \beta, \gamma)
  \;\le\;
  \sqrt{LO(G; \alpha_1, \beta_1, \gamma_1)\,
        LO(G; \alpha_2, \beta_2, \gamma_2)},
  \]
  which is the inequality.

  The Cauchy--Schwarz inequality is an equality if and only if the
  sequences $(a_e)$ and $(b_e)$ are proportional, i.e.\ there exists
  a constant $\lambda \ge 0$ such that $b_e = \lambda\, a_e$ for
  every $e \in E(G)$, equivalently
  $\psi_e^{(2)}/\psi_e^{(1)} = \lambda^{2}$ is the same constant on
  every edge.
  \end{proof}

  \begin{proposition}\label{thm:upper_Delta}
  Let $G$ be a graph of size $m$ and maximum degree $\Delta$. For
  $\alpha, \beta \ge 0$ and $\gamma \ge 0$,
  \begin{equation}\label{eq:upper_Delta}
  \Wuzi \;\le\; m \cdot \Delta^{2\alpha} \cdot (2\Delta)^{\beta}
  \cdot \exp\!\left(\gamma\,\frac{\Delta - 1}{\Delta + 1}\right).
  \end{equation}
  \end{proposition}

  \begin{proof}
  Fix an edge $uv \in E(G)$. By the definition of $\Delta$ we have
  $1 \le d_u, d_v \le \Delta$. We bound the three multiplicative
  factors of $\psi(d_u, d_v; \alpha, \beta, \gamma)$ separately.

  \emph{First factor.} Because $d_u, d_v \le \Delta$ we obtain
  $d_u d_v \le \Delta^{2}$. Since $\alpha \ge 0$ and the map
  $t \mapsto t^{\alpha}$ is non-decreasing on $\mathbb{R}_{> 0}$,
  $(d_u d_v)^{\alpha} \le \Delta^{2\alpha}$.

  \emph{Second factor.} Similarly $d_u + d_v \le 2\Delta$, and the
  map $t \mapsto t^{\beta}$ is non-decreasing for $\beta \ge 0$, so
  $(d_u + d_v)^{\beta} \le (2\Delta)^{\beta}$.

  \emph{Third factor.} The normalised imbalance
  $\mathrm{imb}(x, y) = |x - y|/(x + y)$ is maximised over the
  admissible set $\{(i, j) : 1 \le i \le j \le \Delta\}$ at the
  pair $(1, \Delta)$, where
  $\mathrm{imb}(1, \Delta) = (\Delta - 1)/(\Delta + 1)$. Hence
  $|d_u - d_v|/(d_u + d_v) \le (\Delta - 1)/(\Delta + 1)$ for every
  edge. Because $\gamma \ge 0$ and $\exp$ is non-decreasing,
  $\exp(\gamma\,|d_u - d_v|/(d_u + d_v)) \le \exp(\gamma(\Delta - 1)/(\Delta + 1))$.
  Multiplying the three factor-wise bounds and summing over the
  $m$ edges of $G$ yields~\eqref{eq:upper_Delta}. We note that the
  three per-factor maxima are not simultaneously attainable for
  $\gamma > 0$, so the bound, while valid, is not sharp in that range.
  \end{proof}

  \begin{proposition}\label{thm:lower_delta}
  Let $G$ be a graph of size $m$ and minimum degree $\delta$. For
  $\alpha, \beta \ge 0$ and $\gamma \ge 0$,
  \begin{equation}\label{eq:lower_delta}
  \Wuzi \;\ge\; m \cdot \delta^{2\alpha} \cdot (2\delta)^{\beta},
  \end{equation}
  with equality if and only if $G$ is $\delta$-regular.
  \end{proposition}

  \begin{proof}
  Fix an arbitrary edge $uv \in E(G)$. By the definition of
  $\delta = \delta(G)$ we have $d_u, d_v \ge \delta$, and we bound
  the three multiplicative factors of
  $\psi(d_u, d_v; \alpha, \beta, \gamma)$ separately. Because
  $d_u, d_v \ge \delta$ and $\alpha \ge 0$, $(d_u d_v)^{\alpha} \ge
  \delta^{2\alpha}$; similarly $(d_u + d_v)^{\beta} \ge (2\delta)^{\beta}$;
  and since $\gamma \ge 0$ and the imbalance is non-negative,
  $\exp(\gamma\,|d_u - d_v|/(d_u + d_v)) \ge 1$. Multiplying the three
  lower bounds gives $\psi(d_u, d_v; \alpha, \beta, \gamma) \ge
  \delta^{2\alpha}(2\delta)^{\beta}$ for every edge, and summing over
  the $m$ edges yields~\eqref{eq:lower_delta}.

  For the equality case, the first two per-edge inequalities are
  equalities if and only if $d_u = d_v = \delta$, and under this
  condition the third factor is automatically $1$. Hence equality
  holds throughout if and only if $G$ is $\delta$-regular.
  \end{proof}

  \begin{theorem}\label{thm:m1_bound}
  Let $\alpha = \gamma = 0$ and $\beta \ge 1$. Then
  \[
  \Wuzi
  \;=\; \sum_e (d_u + d_v)^{\beta}
  \;\ge\; m^{1 - \beta}\, M_1(G)^{\beta},
  \]
  where $M_1(G) = \sum_{v \in V} d_v^2 = \sum_{uv \in E}(d_u + d_v)$.
  At $\beta = 1$ both sides equal $M_1(G)$ identically; for
  $\beta > 1$ equality holds if and only if $d_u + d_v$ is constant
  on every edge of $G$.
  \end{theorem}

  \begin{proof}
  Setting $\alpha = \gamma = 0$ in Definition~\ref{def:wuzi}
  collapses the edge-contribution function to
  $\psi(d_u, d_v; 0, \beta, 0) = (d_u + d_v)^{\beta}$, so
  $LO(G; 0, \beta, 0) = \sum_{uv \in E(G)} (d_u + d_v)^{\beta}$.
  The function $\varphi(t) = t^{\beta}$ satisfies
  $\varphi''(t) = \beta(\beta - 1)\, t^{\beta - 2} \ge 0$ for
  $\beta \ge 1$; hence $\varphi$ is convex on
  $\mathbb{R}_{> 0}$, with strict convexity for $\beta > 1$. Apply
  Jensen's inequality (convex case) to the $m = |E(G)|$ positive
  values $\{d_u + d_v : uv \in E(G)\}$ with uniform weights $1/m$:
  \[
  \left(\frac{M_1(G)}{m}\right)^{\!\beta}
  = \varphi\!\left(\frac{1}{m}\sum_{uv} (d_u + d_v)\right)
  \;\le\;
  \frac{1}{m}\sum_{uv} (d_u + d_v)^{\beta}
  = \frac{LO(G; 0, \beta, 0)}{m},
  \]
  using $\sum_{uv \in E(G)} (d_u + d_v) = M_1(G)$. Multiplying both
  sides by $m > 0$ yields the bound. At $\beta = 1$, $\varphi$ is
  affine and Jensen is the identity $LO(G; 0, 1, 0) = M_1(G)$; for
  $\beta > 1$ strict convexity gives equality if and only if
  $d_u + d_v$ is constant over every edge.
  \end{proof}

  \begin{theorem}\label{thm:m1_concave}
  Let $\alpha = \gamma = 0$ and $0 < \beta < 1$. Then
  $\Wuzi = \sum_e (d_u + d_v)^{\beta} \le m^{1 - \beta}\, M_1(G)^{\beta}$,
  with equality if and only if $d_u + d_v$ is constant on every
  edge of $G$.
  \end{theorem}

  \begin{proof}
  The function $t \mapsto t^{\beta}$ is strictly concave on
  $\mathbb{R}_{> 0}$ for $\beta \in (0, 1)$. Jensen's inequality in
  the concave case applied to the values $\{d_u + d_v : uv \in E(G)\}$
  gives $\frac{1}{m} \sum_e (d_u + d_v)^{\beta} \le ( M_1(G) / m )^{\beta}$,
  using $M_1(G) = \sum_{uv} (d_u + d_v)$. Multiplying by $m$ gives
  the upper bound; strict concavity gives equality if and only if
  $d_u + d_v$ is constant on every edge.
  \end{proof}

  \begin{theorem}\label{thm:m2_alpha}
  Let $\beta = \gamma = 0$ and $\alpha \ge 1$. Then
  \[
  \Wuzi
  \;=\; \sum_{uv \in E(G)} (d_u d_v)^{\alpha}
  \;=\; R_{\alpha}(G)
  \;\ge\; m^{\,1 - \alpha}\, M_2(G)^{\alpha},
  \]
  where $M_2(G) = \sum_{uv \in E(G)} d_u d_v$ is the second Zagreb
  index and $R_{\alpha}(G) = \sum_{uv \in E(G)} (d_u d_v)^{\alpha}$
  is the general Randi\'c index. At $\alpha = 1$ both sides equal
  $M_2(G)$ identically; for $\alpha > 1$ equality holds if and only
  if the product $d_u d_v$ is constant on every edge of $G$.
  \end{theorem}

  \begin{proof}
  Setting $\beta = \gamma = 0$ collapses the edge contribution to
  $(d_u d_v)^{\alpha}$, so $\Wuzi = R_{\alpha}(G)$. The function
  $\varphi(t) = t^{\alpha}$ has
  $\varphi''(t) = \alpha(\alpha - 1)\, t^{\alpha - 2} \ge 0$ for
  $\alpha \ge 1$ and is convex, strictly for $\alpha > 1$. Jensen
  (convex case) on $\{d_u d_v : uv \in E(G)\}$ with weights $1/m$
  gives $(M_2(G)/m)^{\alpha} \le R_{\alpha}(G)/m$, using
  $\sum_{uv} d_u d_v = M_2(G)$. Multiplying by $m$ yields the bound;
  $\alpha = 1$ is the affine identity, and for $\alpha > 1$ strict
  convexity gives equality if and only if $d_u d_v$ is constant on
  every edge.
  \end{proof}

  \begin{proposition}\label{thm:R_chi}
  Let $\gamma = 0$, and define
  $R_a(G) = \sum_{uv \in E(G)} (d_u d_v)^{a}$ and
  $\chi_b(G) = \sum_{uv \in E(G)} (d_u + d_v)^{b}$ (so $R = R_{-1/2}$
  and $\chi = \chi_{-1/2}$). Then for every graph $G$,
  \begin{equation}\label{eq:R_chi_bound}
  LO(G; \alpha, \beta, 0)
  \;\le\; \sqrt{R_{2\alpha}(G) \cdot \chi_{2\beta}(G)}.
  \end{equation}
  In particular, at $(\alpha, \beta) = (-1/4, -1/4)$ this
  specialises to $LO(G; -1/4, -1/4, 0) \le \sqrt{R(G) \cdot \chi(G)}$.
  \end{proposition}

  \begin{proof}
  Apply Theorem~\ref{thm:cs} with
  $(\alpha_1, \beta_1, \gamma_1) = (2\alpha, 0, 0)$ and
  $(\alpha_2, \beta_2, \gamma_2) = (0, 2\beta, 0)$, whose components
  sum to $(2\alpha, 2\beta, 0)$, meeting the hypotheses since
  $\gamma = 0$. By Definition~\ref{def:wuzi},
  $LO(G; 2\alpha, 0, 0) = R_{2\alpha}(G)$ and
  $LO(G; 0, 2\beta, 0) = \chi_{2\beta}(G)$, so Theorem~\ref{thm:cs}
  gives~\eqref{eq:R_chi_bound}. The equality case is that of
  Theorem~\ref{thm:cs}: the ratio
  $(d_u d_v)^{2\alpha}/(d_u + d_v)^{2\beta}$ is constant over every
  edge.
  \end{proof}

  \begin{proposition}\label{thm:sombor}
  For $\alpha = 0$, $0 < \beta < 1$, and $\gamma = 0$, the Loyola
  quantity $LO(G; 0, \beta, 0) = \chi_{\beta}(G)$ satisfies
  \begin{equation}\label{eq:sombor_bound}
  LO(G; 0, \beta, 0)
  \;\le\; m^{1 - \beta} \cdot \bigl(\sqrt{2}\, SO(G)\bigr)^{\beta},
  \end{equation}
  where $SO(G) = \sum_{uv \in E(G)} \sqrt{d_u^2 + d_v^2}$ is the
  Sombor index.
  \end{proposition}

  \begin{proof}
  For positive reals $x, y$, $x^2 + y^2 \ge (x + y)^2/2$ follows from
  $(x + y)^2 \le 2(x^2 + y^2)$ by the AM--GM inequality
  $xy \le (x^2 + y^2)/2$. Taking square roots with $x = d_u, y = d_v$
  and summing over $E(G)$ gives
  $SO(G) \ge M_1(G)/\sqrt{2}$, i.e.\ $M_1(G) \le \sqrt{2}\, SO(G)$.
  For $0 < \beta < 1$, the concave Jensen bound
  (Theorem~\ref{thm:m1_concave}) gives
  $LO(G; 0, \beta, 0) \le m^{1 - \beta}\, M_1(G)^{\beta}$, and since
  $t \mapsto t^{\beta}$ is non-decreasing for $\beta > 0$,
  substituting $M_1(G) \le \sqrt{2}\, SO(G)$
  yields~\eqref{eq:sombor_bound}.
  \end{proof}

  The next theorem gives a single uniform template for bounding
  $\Wuzi$ both above and below in terms of any classical
  bond-incident-degree index whose edge contribution is strictly
  positive on the admissible degree pairs of the graph.

  \begin{theorem}\label{thm:ratio_bound}
  Let $G \in \mathcal{G}_{\Delta}$ be a simple graph and let
  $I_h(G) = \sum_{uv \in E(G)} h(d_u, d_v)$ be a BID index with
  edge contribution $h(i, j)$ strictly positive on every degree pair
  $(i, j)$ that appears in $G$. Define
  \[
  c_{\min}^{h} = \min_{(i, j)}\,\frac{\psi(i, j; \alpha, \beta,
  \gamma)}{h(i, j)},
  \qquad
  c_{\max}^{h} = \max_{(i, j)}\,\frac{\psi(i, j; \alpha, \beta,
  \gamma)}{h(i, j)},
  \]
  where the extrema range over the admissible degree pairs of $G$.
  Then
  \begin{equation}\label{eq:ratio_bound}
  c_{\min}^{h} \cdot I_h(G)
  \;\le\; \Wuzi
  \;\le\; c_{\max}^{h} \cdot I_h(G),
  \end{equation}
  with equality on either side if and only if every edge of $G$
  attains the corresponding extremum of $\psi/h$.
  \end{theorem}

  \begin{proof}
  Fix an edge $uv \in E(G)$ and define
  $r_{uv} := \psi(d_u, d_v; \alpha, \beta, \gamma)/h(d_u, d_v) > 0$.
  By definition $c_{\min}^{h} \le r_{uv} \le c_{\max}^{h}$, so
  multiplying by $h(d_u, d_v) > 0$ gives the per-edge sandwich
  $c_{\min}^{h} h(d_u, d_v) \le \psi(d_u, d_v; \alpha, \beta, \gamma)
  \le c_{\max}^{h} h(d_u, d_v)$. Summing over $E(G)$ and using
  $\sum_{uv} h(d_u, d_v) = I_h(G)$ and $\sum_{uv} \psi = \Wuzi$
  gives~\eqref{eq:ratio_bound}; the equality conditions are that
  every edge attains the corresponding extremum of $\psi/h$.
  \end{proof}

  \begin{proposition}\label{prop:gamma_eq_0_or_inf}
  At $\gamma = 0$, $LO(G; \alpha, \beta, 0)$ coincides with the
  two-parameter $(\alpha, \beta)$ degree-based index
  $\sum_{uv \in E(G)} (d_u d_v)^{\alpha}(d_u + d_v)^{\beta}$.
  Let $E^{*}(G)$ denote the set of edges realising
  $\mathrm{imb}_{\max}(G) = \max_{uv} \mathrm{imb}(d_u, d_v)$
  and $E_{*}(G)$ the set realising
  $\mathrm{imb}_{\min}(G) = \min_{uv} \mathrm{imb}(d_u, d_v)$.
  Then as $\gamma \to +\infty$,
  \[
  \frac{\Wuzi}{e^{\gamma\, \mathrm{imb}_{\max}(G)}}
  \;\longrightarrow\;
  \sum_{uv \in E^{*}(G)} (d_u d_v)^{\alpha}(d_u + d_v)^{\beta},
  \]
  and as $\gamma \to -\infty$,
  $\Wuzi / e^{\gamma\, \mathrm{imb}_{\min}(G)} \longrightarrow
  \sum_{uv \in E_{*}(G)} (d_u d_v)^{\alpha}(d_u + d_v)^{\beta}$.
  \end{proposition}

  \begin{proof}
  At $\gamma = 0$ the exponential factor is $1$ on every edge.
  For the asymptotics, write $M(G) := \mathrm{imb}_{\max}(G)$ and
  $\Delta_{uv} := M(G) - \mathrm{imb}(d_u, d_v) \ge 0$, so that
  $\Wuzi / e^{\gamma M(G)} = \sum_{uv} (d_u d_v)^{\alpha}(d_u +
  d_v)^{\beta} e^{-\gamma \Delta_{uv}}$. For $uv \in E^{*}(G)$,
  $\Delta_{uv} = 0$ and the summand is $\gamma$-independent; for
  $uv \notin E^{*}(G)$, $\Delta_{uv} > 0$ and $e^{-\gamma \Delta_{uv}}
  \to 0$ as $\gamma \to +\infty$. Since $E(G)$ is finite the limit
  is termwise, giving the identity. The $\gamma \to -\infty$ case is
  dual.
  \end{proof}

% =====================================================================
\section{Descriptor-space saturation}\label{sec:saturation}
% =====================================================================

Section~\ref{sec:intro} observed that on molecular graphs every BID
index is a linear functional of the edge-degree-pair counts. We now
make this precise for the Loyola family, and show that the resulting
bound is not merely a containment but an exact equality: the three
continuous parameters realise the whole available descriptor space,
and nothing beyond it.

Fix $\Delta_0 \in \mathbb{N}$ and let
$\mathcal{G} = \{G_1, \dots, G_N\}$ be graphs with
$\Delta(G_k) \le \Delta_0$. For $1 \le i \le j \le \Delta_0$ write
$m_{ij}(G) = |\{uv \in E(G) : \{d_u, d_v\} = \{i, j\}\}|$ and let
$M = [\,m_{ij}(G_k)\,] \in \mathbb{R}^{N \times D}$ with
$D = \binom{\Delta_0 + 1}{2} = \tfrac{\Delta_0(\Delta_0 + 1)}{2}$.

\begin{theorem}\label{thm:saturation}
For every $(\alpha, \beta, \gamma) \in \mathbb{R}^3$,
\[
\bigl(LO(G_k; \alpha, \beta, \gamma)\bigr)_{k=1}^{N}
= M\,\psi(\alpha, \beta, \gamma),
\qquad
\psi(\alpha, \beta, \gamma)
= \bigl(\psi(i, j; \alpha, \beta, \gamma)\bigr)_{1 \le i \le j \le \Delta_0},
\]
and moreover
\[
\operatorname{span}\bigl\{(LO(G_k; \alpha, \beta, \gamma))_k :
(\alpha, \beta, \gamma) \in \mathbb{R}^3\bigr\}
= \operatorname{col}(M),
\qquad
\dim = \operatorname{rank}(M) \le \tfrac{\Delta_0(\Delta_0 + 1)}{2}.
\]
In particular, on hydrogen-suppressed molecular graphs
$(\Delta_0 = 4)$ the entire three-parameter family spans a space of
dimension exactly $\operatorname{rank}(M) \le 10$, attained but never
exceeded.
\end{theorem}

\begin{proof}
Since $\psi(d_u, d_v; \alpha, \beta, \gamma)$ depends only on the
unordered pair $\{d_u, d_v\}$, partitioning $E(G)$ by degree-pair
class gives
$LO(G; \alpha, \beta, \gamma)
= \sum_{1 \le i \le j \le \Delta_0} \psi(i, j; \alpha, \beta, \gamma)\,
m_{ij}(G)$, i.e.\ the stated vector identity. Writing the edge weight
in exponential form,
\[
\psi(i, j; \alpha, \beta, \gamma)
= \exp\!\bigl(\alpha \ln(ij) + \beta \ln(i + j)
+ \gamma\, \mathrm{imb}(i, j)\bigr)
= \exp\langle \theta, v_{ij} \rangle,
\]
with $\theta = (\alpha, \beta, \gamma)$ and frequency vector
$v_{ij} = (\ln(ij), \ln(i + j), \mathrm{imb}(i, j))$. The $D$
admissible pairs yield $D$ \emph{distinct} frequency vectors: the
sum $i + j$ and product $ij$ determine $\{i, j\}$, so distinct
classes already differ in their first two coordinates. Exponentials
$\theta \mapsto \exp\langle \theta, v \rangle$ with distinct
frequencies are linearly independent as functions of $\theta$; hence
the only covector annihilating every $\psi(\theta)$ is $0$, so
$\{\psi(\theta) : \theta \in \mathbb{R}^3\}$ spans all of
$\mathbb{R}^{D}$. Therefore
$\operatorname{span}\{M\psi(\theta)\} = M\,\mathbb{R}^{D}
= \operatorname{col}(M)$, of dimension $\operatorname{rank}(M) \le D$.
For $\Delta_0 = 4$ the ten frequency vectors are distinct and the
realised rank equals $10$.
\end{proof}

Theorem~\ref{thm:saturation} sharpens the descriptor-space remark of
Section~\ref{sec:intro} from a containment to an exact span. No
amount of continuous parameter freedom in a BID index escapes the
finite basis of edge-degree-pair counts: the Loyola family is a
smooth three-parameter surface filling a fixed finite-dimensional
space, not a route out of it. This is the structural reason the
chemometric study of Section~\ref{sec:sensitivity} finds the family
in the same correlation regime as the classical indices.

% =====================================================================
\section{A \texorpdfstring{$\gamma$}{gamma}-dependent extremal characterization}\label{sec:extremal}
% =====================================================================

While Theorem~\ref{thm:saturation} shows the imbalance parameter
$\gamma$ adds no descriptor \emph{dimension}, $\gamma$ is decisive
for extremal problems: it controls which tree maximises the index.
Throughout, $\mathcal{T}_n$ denotes the set of trees of order $n$.

\begin{theorem}\label{thm:switch}
Let $n \ge 4$ and $\alpha, \beta \in \mathbb{R}$.
\begin{enumerate}
\item There is a threshold $\gamma_0 = \gamma_0(\alpha, \beta, n)$
such that for every $\gamma > \gamma_0$ the star $S_n$ is the unique
maximiser of $LO(\cdot\,; \alpha, \beta, \gamma)$ over
$\mathcal{T}_n$; one may take
\[
\gamma_0 = \max\!\Bigl(0,\ \tfrac{n(n-1)}{2}
\ln \tfrac{m_n}{(n-1)^{\alpha} n^{\beta}}\Bigr),
\qquad
m_n = \max_{1 \le i \le j \le n-2} (ij)^{\alpha}(i+j)^{\beta}.
\]
\item For $(\alpha, \beta) = (-\tfrac12, 0)$,
$LO(\cdot\,; -\tfrac12, 0, 0) = R$ is the Randi\'c index, whose
unique maximiser over $\mathcal{T}_n$ is the path
$P_n$~\cite{BollobasErdos1998}; by continuity $P_n$ remains the
maximiser for $\gamma$ in a right-neighbourhood of $0$.
\end{enumerate}
Hence the maximising tree depends on $\gamma$ --- the path for small
$\gamma$, the star for large $\gamma$ --- so $\gamma$ is not
redundant for extremal problems.
\end{theorem}

\begin{proof}
(1) Every edge of $S_n$ is the pair $(1, n-1)$ with imbalance
$\tfrac{n-2}{n}$. A non-star tree $T$ has $\Delta(T) \le n-2$, so
each of its edges has imbalance at most $\tfrac{n-3}{n-1}$, and
\[
\frac{n-2}{n} - \frac{n-3}{n-1} = \frac{2}{n(n-1)} > 0;
\]
thus every star edge strictly exceeds every non-star edge in
imbalance. With $m_n$ as above (finite and positive) and since a
tree has $n-1$ edges, for $\gamma \ge 0$,
\[
LO(T) \le (n-1)\, m_n\, e^{\gamma \frac{n-3}{n-1}},
\qquad
LO(S_n) = (n-1)(n-1)^{\alpha} n^{\beta}\, e^{\gamma \frac{n-2}{n}},
\]
whence
\[
\frac{LO(S_n)}{LO(T)}
\ge \frac{(n-1)^{\alpha} n^{\beta}}{m_n}\,
\exp\!\Bigl(\gamma \cdot \tfrac{2}{n(n-1)}\Bigr)
\xrightarrow[\gamma \to \infty]{} \infty.
\]
As $\mathcal{T}_n$ is finite, the displayed $\gamma_0$ makes
$LO(S_n) > LO(T)$ for every non-star $T$, with strict inequality
giving uniqueness. (2) At $\gamma = 0$,
$LO(\cdot\,; -\tfrac12, 0, 0) = \sum_{uv} (d_u d_v)^{-1/2} = R$,
whose maximum over $\mathcal{T}_n$ is attained uniquely by $P_n$.
Each map $\gamma \mapsto LO(T; -\tfrac12, 0, \gamma)$ is continuous
and $\mathcal{T}_n$ is finite, so $P_n$ remains the strict maximiser
on a right-neighbourhood of $\gamma = 0$.
\end{proof}

A direct computation over all trees of order $n$ confirms the
transition (Table~\ref{tab:switch}): the maximiser is the path at
$\gamma = 0$ and the star for large $\gamma$, passing through
intermediate (broom-like) trees in between; characterising the
intermediate maximisers is an open problem.

\begin{table}[H]
\centering
\caption{Maximiser of $LO(\cdot\,; -1/2, 0, \gamma)$ over all trees of order $n$, computed exactly over the full tree set.}
\label{tab:switch}
\small
\begin{tabular}{lcccc}
\toprule
$\gamma$ & $n = 7$ & $n = 8$ & $n = 9$ & $n = 10$\\
\midrule
$0$   & path  & path  & path  & path\\
$0.5$ & other & other & other & other\\
$1$   & star  & other & other & other\\
$2$   & star  & star  & star  & star\\
$5$   & star  & star  & star  & star\\
\bottomrule
\end{tabular}
\end{table}

% =====================================================================
\section{Sensitivity analysis}\label{sec:sensitivity}
% =====================================================================

We benchmark the Loyola family against the classical indices it
reduces to, on three standard testbeds, following the chemometric
benchmark of the diminished Sombor
study~\cite{MovahediGutmanRedzepovicFurtula2026}: the $18$ octane
isomers (five NIST physicochemical properties), the $106$
non-isomorphic trees of order $10$, and the $75$ decane isomers. The comparison set
is restricted to the exact reductions of Section~\ref{sec:intro} ---
$M_1, M_2, HM, {}^{m}\!M_2, R, \chi, H/2, ISI, GA, AG$ --- together
with the two \emph{pre-specified} pure-$\gamma$ points $LO(0, 0, 1)$
and $LO(0, 0, 2)$, fixed in advance. No parameter is tuned against
the property data; the edge count $m = LO(0, 0, 0)$ is omitted as it
is constant on the octanes.

\subsection{Property correlations on the octane isomers}\label{sec:octanes}

Table~\ref{tab:lo_octane} reports the Pearson correlation $r$ between
each index and each of the five physicochemical properties of the
$18$ octane isomers (boiling point $T_B$, enthalpy of formation
$\Delta H_f$, enthalpy of vaporization $\Delta H_{\mathrm{vap}}$,
entropy $S$, acentric factor $\omega$; property data from the NIST
WebBook).

\begin{table}[H]
\centering
\caption{Pearson correlation $r$ between the Loyola reductions and the pre-specified pure-$\gamma$ Loyola points and the five physicochemical properties of the $18$ octane isomers (NIST WebBook). Bold flags the best $|r|$ per column. The edge count $m = LO(0,0,0)$ is omitted (constant on octanes).}
\label{tab:lo_octane}
\small
\begin{tabular}{lrrrrr}
\toprule
Index & $T_B$ & $\Delta H_f$ & $\Delta H_{\mathrm{vap}}$ & $S$ & $\omega$\\
\midrule
$M_1 = LO(0,1,0)$ & $-0.718$ & $-0.762$ & $-0.912$ & $-0.973$ & $-0.970$\\
$M_2 = LO(1,0,0)$ & $-0.503$ & $-0.553$ & $-0.775$ & $-0.942$ & $\mathbf{-0.987}$\\
$HM = LO(0,2,0)$ & $-0.655$ & $-0.712$ & $-0.873$ & $\mathbf{-0.975}$ & $-0.982$\\
${}^{m}\!M_2 = LO(-1,0,0)$ & $\mathbf{+0.861}$ & $\mathbf{+0.888}$ & $+0.942$ & $+0.866$ & $+0.799$\\
$R = LO(-1/2,0,0)$ & $+0.821$ & $+0.848$ & $+0.956$ & $+0.935$ & $+0.901$\\
$\chi = LO(0,-1/2,0)$ & $+0.801$ & $+0.829$ & $+0.953$ & $+0.950$ & $+0.927$\\
$H/2 = LO(0,-1,0)$ & $+0.822$ & $+0.846$ & $+0.959$ & $+0.935$ & $+0.904$\\
$ISI = LO(1,-1,0)$ & $-0.037$ & $-0.047$ & $-0.363$ & $-0.632$ & $-0.760$\\
$GA = 2\,LO(1/2,-1,0)$ & $+0.823$ & $+0.856$ & $+0.959$ & $+0.942$ & $+0.909$\\
$AG = \tfrac12 LO(-1/2,1,0)$ & $-0.817$ & $-0.856$ & $-0.951$ & $-0.941$ & $-0.903$\\
$LO(0,0,1)$ & $-0.828$ & $-0.831$ & $-0.966$ & $-0.931$ & $-0.930$\\
$LO(0,0,2)$ & $-0.830$ & $-0.844$ & $\mathbf{-0.967}$ & $-0.939$ & $-0.926$\\
\bottomrule
\end{tabular}
\end{table}

The best correlation per property is split among the reductions and
the pure-$\gamma$ points: the modified second Zagreb index
$LO(-1, 0, 0)$ is strongest on $T_B$ and $\Delta H_f$, the
hyper-Zagreb index $LO(0, 2, 0)$ on $S$, the second Zagreb index
$LO(1, 0, 0)$ on $\omega$, and the pre-specified pure-$\gamma$ point
$LO(0, 0, 2)$ on $\Delta H_{\mathrm{vap}}$ ($|r| = 0.967$). At the
exact-reduction triples the Loyola correlations coincide with the
corresponding classical indices to floating-point precision, as they
must. All per-property margins are within the bootstrap confidence
half-width at $n = 18$ (a percentile bootstrap with
$B = 10\,000$ resamples gives half-widths of order $0.03$ on the
headline cells), so no index --- classical reduction or Loyola point
--- is significantly separated from its competitors on this small
set; the value of $LO(0, 0, 2)$ on $\Delta H_{\mathrm{vap}}$ is a
pre-specified, not a tuned, result. A leave-one-out check across the
$18$ isomers shifts every headline $|r|$ by at most $0.011$,
confirming the correlations are not driven by a single outlying
isomer.

Figure~\ref{fig:octanes} shows the signed-$r$ heatmap of the
compared indices against the five properties.

\begin{figure}[H]
\centering
\includegraphics[width=0.9\linewidth]{fig_lo_prediction_correlation_heatmap.pdf}
\caption{Signed Pearson $r$ between the Loyola reductions and pre-specified pure-$\gamma$ points and the five octane physicochemical properties on the $18$ isomers. The best absolute correlation per property is marked.}
\label{fig:octanes}
\end{figure}

On the octanes the pure-$\gamma$ points are highly collinear with
the classical reductions (maximum $|r|$ with a reduction equal to
$0.983$ for $LO(0, 0, 1)$ and $0.993$ for $LO(0, 0, 2)$), consistent
with Theorem~\ref{thm:saturation}: the parameter $\gamma$ moves the
descriptor within the fixed BID span rather than outside it.

\subsection{Degeneracy on trees of order 10}\label{sec:degeneracy}

Degeneracy measures an index's ability to discriminate
non-isomorphic graphs~\cite{Konstantinova1996Degeneracy}; lower is
better. On the $106$ non-isomorphic
trees of order $10$ we use the distinct-value convention
\[
\mathrm{Degeneracy}(\%) = 100\Bigl(1 - \tfrac{k}{N}\Bigr),
\qquad k = \#\{\text{distinct index values}\}, \quad N = 106,
\]
i.e.\ the fraction of the $N$ structures that fail to receive a
unique value. The pure-$\gamma$ points $LO(0, 0, 1)$ and
$LO(0, 0, 2)$ attain degeneracy $25.5\%$ ($79$ distinct values),
among the lowest of all compared indices and tied with $R$, $GA$,
and $AG$; the first Zagreb index $LO(0, 1, 0)$ ($83.0\%$) and the
second Zagreb index $LO(1, 0, 0)$ ($67.0\%$) are the most degenerate
(Figure~\ref{fig:degeneracy}).

\begin{figure}[H]
\centering
\includegraphics[width=0.9\linewidth]{fig_lo_degeneracy_order10_trees.pdf}
\caption{Degeneracy percentage of the compared indices on the $106$ non-isomorphic trees of order $10$ (lower is better).}
\label{fig:degeneracy}
\end{figure}

\subsection{Structure sensitivity on decane isomers}\label{sec:ss}

The structure-sensitivity criterion for descriptor quality was
introduced by Furtula, Gutman and
Dehmer~\cite{FurtulaGutmanDehmer2013SS}, with a fingerprint-based
variant due to Rak\'i\'c and Furtula~\cite{RakicFurtula2019SS}. As a
simple, self-contained proxy used only for within-study comparison,
we report over the $75$ decane isomers
the coefficient of variation $SS = \sigma/\mu$, the
relative range $Abr = (\max - \min)/\mu$, and their ratio
$SA = SS/Abr$ (Figure~\ref{fig:ss}). The pure-$\gamma$ point
$LO(0, 0, 2)$ has the highest structure sensitivity ($SS = 0.165$)
and the largest relative range among all compared indices, with
$LO(0, 0, 1)$ also high; the imbalance parameter's contribution
surfaces here, in structural responsiveness, rather than in added
linear correlation.

\begin{figure}[H]
\centering
\includegraphics[width=0.95\linewidth]{fig_lo_structure_sensitivity_decanes.pdf}
\caption{Structure-sensitivity measures $SS = \sigma/\mu$, $Abr = (\max-\min)/\mu$, and $SA = SS/Abr$ of the compared indices on the $75$ decane isomers.}
\label{fig:ss}
\end{figure}

% =====================================================================
\section{Conclusion and further directions}\label{sec:conclusion}
% =====================================================================

The Loyola family is a three-parameter degree-based index that
recovers many classical indices as exact reductions and lies
inside the same finite-dimensional descriptor space as the
established indices. We have given closed-form values on standard
graph classes and several sharp upper and lower bounds: a
bracket bound tight on regular graphs, Cauchy--Schwarz and
Jensen-type bounds, bounds using the basic graph parameters
$n, m, \Delta, \delta$, bounds via
the generalised Randi\'c, sum-connectivity, and Sombor indices, and a uniform ratio bound against any
classical index with positive edge contribution. We have further
proved that the family's descriptor span is exactly the column space
of the edge-degree-pair count matrix (Theorem~\ref{thm:saturation}),
and that the imbalance parameter, though adding no descriptor
dimension, controls the extremal tree, switching the maximiser from
the path to the star (Theorem~\ref{thm:switch}).

Chemometrically, a sensitivity benchmark on $18$ octane isomers, $106$
non-isomorphic trees of order $10$, and $75$ decane isomers
places the Loyola family in the same correlation regime as the
classical BID indices. The Loyola family is hence best viewed as a controlled
parametric probe of the finite BID descriptor space.

Ultimately, the biggest open problem would be to propose a new non-redundant topological index that performs better than all of the existing ones in regard to prediction of physiochemical properties. The work in this paper following the Loyola index can be extended by research directions that can be taken by the following open problems:
\begin{enumerate}
\item Sharp two-sided bounds for $\Wuzi$ via popular topological indices, including the
geometric--arithmetic index
$GA$~\cite{VukicevicFurtula2009GA}, the harmonic index $H$, and
the atom-bond-connectivity index $ABC$.


\item Extremal-graph characterisations for further graph classes, including unicyclic and
bicyclic graphs, and a description of the intermediate-$\gamma$ tree maximisers of Theorem~\ref{thm:switch}.

\item Chemometric extension to additional chemical compound classes
(benzene derivatives, alkane homologues beyond octane and
decane) and to compound-class-specific QSPR endpoints.
\end{enumerate}

\acknowledgment{The authors thank colleagues for discussions of
the computational extremal-graph search. Funding sources and institutional
support to be acknowledged in the final version of the
manuscript.}

% References should be ordered alphabetically.
\singlespacing

 \begin{thebibliography}{99}

  \bibitem{AhmadFarooqDas2025Extremal}
  S. Ahmad, R. Farooq, K. C. Das, The general Sombor index of extremal
  trees with a given maximum degree, \textit{MATCH Commun. Math. Comput.
  Chem.\/} \textbf{94} (2025) 825--853.

  \bibitem{Albertson1997}
  M. O. Albertson, The irregularity of a graph, \textit{Ars Combin.\/}
  \textbf{46} (1997) 219--225.

  \bibitem{ArockiarajPaulClement2023Hybrid} M. Arockiaraj, D. Paul, J.
  Clement, S. Tigga, K. Jacob, K. Balasubramanian, Novel molecular hybrid
  geometric-harmonic-Zagreb degree based descriptors and their efficacy in
  QSPR studies of polycyclic aromatic hydrocarbons, \textit{SAR QSAR
  Environ. Res.\/} \textbf{34} (2023) 569--589.

  \bibitem{BarmanDas2026HSO} J. Barman, S. Das, Geometric approach to
  degree-based topological index: Hyperbolic Sombor index, \textit{MATCH
  Commun. Math. Comput. Chem.\/} \textbf{95} (2026) 63--94.

  \bibitem{BollobasErdos1998}
  B. Bollobás, P. Erdős, Graphs of extremal weights, \textit{Ars Combin.\/}
  \textbf{50} (1998) 225--233.

  \bibitem{DasCevikCangulShang2021Sombor} K. C. Das, A. S. \c{C}evik, I. N.
  Cangul, Y. Shang, On Sombor index, \textit{Symmetry\/} \textbf{13}
  (2021) \#140.

  \bibitem{DasGutman2022SomborTrees} K. C. Das, I. Gutman, On Sombor index
  of trees, \textit{Appl. Math. Comput.\/} \textbf{412} (2022) \#126575.

  \bibitem{FurtulaGutman2015F} B. Furtula, I. Gutman, A forgotten
  topological index, \textit{J. Math. Chem.\/} \textbf{53} (2015)
  1184--1190.

  \bibitem{FurtulaGutmanDehmer2013SS} B. Furtula, I. Gutman, M. Dehmer, On
  structure-sensitivity of degree-based topological indices, \textit{Appl.
  Math. Comput.\/} \textbf{219} (2013) 8973--8978.

  \bibitem{GujiWali2025Lanzhou} R. Guji, M. Wali, The Sombor index
  (coindex) and Lanzhou index (coindex) of some graphs, \textit{Axioms\/}
  \textbf{14} (2025) 164.

  \bibitem{GutmanToganYurttasCevikCangul2018Sigma} I. Gutman, M. Togan, A.
  Yurttaş, A. S. Çevik, I. N. Cangul, Inverse problem for sigma index,
  \textit{MATCH Commun. Math. Comput. Chem.\/} \textbf{79} (2018) 491--508.

  \bibitem{Gutman2021Sombor} I. Gutman, Geometric approach to degree-based
  topological indices: Sombor indices, \textit{MATCH Commun. Math. Comput.
  Chem.\/} \textbf{86} (2021) 11--16.

  \bibitem{Gutman2021SomborBasic} I. Gutman, Some basic properties of
  Sombor indices, \textit{Open J. Discr. Appl. Math.\/} \textbf{4} (2021)
  1--3.

  \bibitem{GutmanFurtulaOz2024Elliptic} I. Gutman, B. Furtula, M. S. Oz,
  Geometric approach to vertex-degree-based topological indices: Elliptic
  Sombor index, theory and application, \textit{Int. J. Quantum Chem.\/}
  \textbf{124} (2024) \#e27346.

  \bibitem{Gutman2026SomborSurvey} I. Gutman, Survey of Sombor-type
  degree-based topological indices, \textit{MATCH Commun. Math. Comput.
  Chem.\/} \textbf{96} (2026) 819--837.

  \bibitem{KalaamGreeniArockiaraj2024} A. R. A. Kalaam, A. B. Greeni, M.
  Arockiaraj, Modified reverse degree descriptors for combined topological
  and entropy characterizations of 2D metal organic frameworks:
  applications in graph energy prediction, \textit{Front. Chem.\/}
  \textbf{12} (2024) \#1470231.

  \bibitem{Konstantinova1996Degeneracy} E. V. Konstantinova, The
  discrimination ability of some topological and information distance indices
  for graphs of unbranched hexagonal systems, \textit{J. Chem. Inf. Comput.
  Sci.\/} \textbf{36} (1996) 54--57.

  \bibitem{KosariDehgardiKhan2023Lower} S. Kosari, N. Dehgardi, A. Khan,
  Lower bound on the KG-Sombor index, \textit{Commun. Comb. Optim.\/}
  \textbf{8} (2023) 751--757.

  \bibitem{MendezBermudezAguilarSanchez2022MeanSombor} J. A.
  Méndez-Bermúdez, R. Aguilar-Sánchez, E. D. Molina, J. M. Rodríguez, Mean
  Sombor index, \textit{Discr. Math. Lett.\/} \textbf{9} (2022) 18--25.

  \bibitem{MilovanovicMilovanovicMatejic2021Sombor} I. Milovanović, E.
  Milovanović, M. Matejić, On some mathematical properties of Sombor
  indices, \textit{Bull. Int. Math. Virtual Inst.\/} \textbf{11} (2021)
  341--353.

  \bibitem{Movahedi2025DSO} F. Movahedi, Diminished Sombor index and its
  relationship with topological indices, \textit{Filomat\/} \textbf{39}
  (2025) 10519--10532.

  \bibitem{MovahediGutmanRedzepovicFurtula2026} F. Movahedi, I. Gutman, I.
  Redžepović, B. Furtula, Diminished Sombor index, \textit{MATCH Commun.
  Math. Comput. Chem.\/} \textbf{95} (2026) 141--162.

  \bibitem{Rada2026LinearVDB} J. Rada, A finite-dimensional linear
  framework for vertex-degree-based topological indices, \textit{MATCH
  Commun. Math. Comput. Chem.\/} \textbf{96} (2026) 841--863.

  \bibitem{RakicFurtula2019SS} M. Raki\'c, B. Furtula, A novel method for
  measuring the structure sensitivity of molecular descriptors, \textit{J.
  Chemometrics\/} \textbf{33} (2019) e3138.

  \bibitem{Shang2022Simplicial}
  Y. Shang, Sombor index and degree-related properties of simplicial
  networks, \textit{Appl. Math. Comput.\/} \textbf{419} (2022) 126881.

  \bibitem{VukicevicFurtula2009GA} D. Vukičević, B. Furtula, Topological
  index based on the ratios of geometrical and arithmetical means of
  end-vertex degrees of edges, \textit{J. Math. Chem.\/} \textbf{46} (2009)
  1369--1376.

  \bibitem{WangMaoLiFurtula2022SomborRelations} Z. Wang, Y. Mao, Y. Li, B.
  Furtula, On relations between Sombor and other degree-based indices,
  \textit{J. Appl. Math. Comput.\/} \textbf{68} (2022) 1--17.

  \bibitem{ZhouTrinajstic2009sumconn}
  B. Zhou, N. Trinajstić, On a novel connectivity index, \textit{J. Math.
  Chem.\/} \textbf{46} (2009) 1252--1270.

  \bibitem{ZuoRatherImranAli2022Modified}
  X. Zuo, B. A. Rather, M. Imran, A. Ali, On some topological indices
  defined via the modified Sombor matrix, \textit{Molecules\/} \textbf{27}
  (2022) 6772.

  \end{thebibliography}

\end{document}
